\chapter{Flat base change for the pushforward of a quasi-coherent sheaf ($i=0$)}
\label{chap:Cohomology_FlatBaseChange}
% SOURCE: [Stacks Project], Cohomology of Schemes, Section "Cohomology and
% base change, I", and Tag 02KH (read from references/stacks-coherent.tex and
% references/stacks-coherent.md).

\section{Setup and motivation}

This chapter establishes the $i=0$ (i.e.\ \(H^0\) / direct-image) case of flat
base change: the formation of the pushforward \(f_*\mathcal{F}\) of a
quasi-coherent sheaf commutes with flat base change on the target. It is the
foundational, fully self-contained slice of the cohomology-and-base-change
package that ultimately underlies relative representability statements; the
higher derived cases \(R^i f_*\) for \(i \geq 1\) are treated separately and are
out of scope here.

Throughout, \(f : X \to S\) is a morphism of schemes that is quasi-compact and
quasi-separated, and \(\mathcal{F}\) is a quasi-coherent \(\mathcal{O}_X\)-module.
A base change of \(f\) along an arbitrary morphism \(g : S' \to S\) is recorded
by the cartesian square obtained from the fibre product
\(X' = X \times_S S'\):
\[
  \begin{array}{ccc}
    X' & \xrightarrow{\;g'\;} & X \\[2pt]
    {\scriptstyle f'}\big\downarrow & & \big\downarrow{\scriptstyle f} \\[2pt]
    S' & \xrightarrow{\;g\;} & S
  \end{array}
\]
Here \(f' : X' \to S'\) is the base change of \(f\), and \(g' : X' \to X\) is the
first projection. We write \(\mathcal{F}' = (g')^*\mathcal{F}\) for the pullback
of \(\mathcal{F}\) to \(X'\).

\section{The canonical base-change map}

\begin{definition}\leanok
[Base-change map for the pushforward]
  \label{def:pushforward_base_change_map}
  \lean{AlgebraicGeometry.pushforwardBaseChangeMap}
  % SOURCE: [Stacks Project], Cohomology of Schemes, Section "Cohomology and
  % base change, I", base-change diagram (\ref{equation-base-change-diagram})
  % (read from references/stacks-coherent.tex, L877--904).
  % SOURCE QUOTE: "Let $f : X \to S$ be a morphism of schemes.
  % Let $\mathcal{F}$ be a quasi-coherent sheaf on $X$.
  % Suppose further that $g : S' \to S$ is any morphism of schemes. Denote
  % $X' = X_{S'} = S' \times_S X$ the base change of $X$ and denote
  % $f' : X' \to S'$ the base change of $f$.
  % Also write $g' : X' \to X$ the projection,
  % and set $\mathcal{F}' = (g')^*\mathcal{F}$."
  \textit{Source: Stacks Project, Cohomology of Schemes, \S\,Cohomology and base
  change, I.}
  Let \(f : X \to S\) be a morphism of schemes, \(\mathcal{F}\) a quasi-coherent
  \(\mathcal{O}_X\)-module, and \(g : S' \to S\) an arbitrary morphism, giving the
  cartesian square above with projection \(g' : X' \to X\) and
  \(\mathcal{F}' = (g')^*\mathcal{F}\). The \emph{base-change map for the
  pushforward} is the canonical morphism of \(\mathcal{O}_{S'}\)-modules
  \[
    g^*\bigl(f_*\mathcal{F}\bigr) \longrightarrow f'_*\bigl((g')^*\mathcal{F}\bigr)
    = f'_*\mathcal{F}' .
  \]
  It is adjoint to the morphism
  \(f_*\mathcal{F} \to g_*\, f'_*\mathcal{F}' = f_*\,(g')_*\mathcal{F}'\) obtained
  by applying \(f_*\) to the unit
  \(\mathcal{F} \to (g')_*(g')^*\mathcal{F}\) of the
  \(((g')^*, (g')_*)\)-adjunction, using the commutativity
  \(g \circ f' = f \circ g'\) of the square.
\end{definition}

\section{The affine case}

\subsection{Locality of isomorphisms for \texttt{Scheme.Modules} morphisms}

Before treating the affine computation we record three project-local supplements
to Mathlib. Mathlib provides the per-open criterion
(\(\mathrm{IsIso}\,\varphi \iff \forall U,\ \mathrm{IsIso}(\varphi_U)\)) and a
stalkwise isomorphism criterion for \(\mathrm{TopCat.Sheaf}\)-valued morphisms, but
it does not package the stalk-local or basis-local criterion at the level of
morphisms of \(\mathcal{O}_X\)-modules. The following lemmas bridge that gap; they
are exactly the locality tools used in the first reduction of the affine
base-change lemma below, where one checks the base-change map after restricting to
the affine opens of the base. As bespoke infrastructure they carry no external
citation and stand on their proof sketches.

\begin{lemma}\leanok
[Stalk-local criterion for isomorphisms of \(\mathcal{O}_X\)-modules]
  \label{lem:modules_isIso_iff_stalk}
  \lean{AlgebraicGeometry.Modules.isIso_iff_isIso_stalkFunctor_map}
  Let \(X\) be a scheme and \(\varphi : M \to N\) a morphism of sheaves of
  \(\mathcal{O}_X\)-modules. Then \(\varphi\) is an isomorphism if and only if the
  underlying morphism of abelian-group presheaves induces an isomorphism on the
  stalk at every point \(x \in X\).
\end{lemma}

\begin{proof}


\leanok
  If \(\varphi\) is an isomorphism, then so is its image under the forgetful
  functor to abelian-group presheaves, and the stalk functor at each \(x\),
  being a functor, carries isomorphisms to isomorphisms. Conversely, assume the
  stalk map is an isomorphism at every \(x\). Package the underlying abelian-group
  presheaves of \(M\) and \(N\) as sheaves of abelian groups, and \(\varphi\) as a
  morphism of these sheaves; the stalkwise-isomorphism hypothesis lets us apply the
  \(\mathrm{TopCat.Sheaf}\)-level stalk-local isomorphism criterion to conclude that
  this morphism of sheaves of abelian groups is an isomorphism. Pushing it through
  the forgetful functor to presheaves yields an isomorphism of the underlying
  abelian-group presheaf morphisms, and since the forgetful functor from
  \(\mathcal{O}_X\)-modules to abelian-group presheaves reflects isomorphisms,
  \(\varphi\) itself is an isomorphism.
\end{proof}

\begin{lemma}\leanok
[Basis-local criterion for isomorphisms of \(\mathcal{O}_X\)-modules]
  \label{lem:modules_isIso_of_isBasis}
  \lean{AlgebraicGeometry.Modules.isIso_of_isIso_app_of_isBasis}
  \uses{lem:modules_isIso_iff_stalk}
  Let \(X\) be a scheme, let \(\{B_i\}_{i \in \iota}\) be a family of opens whose
  range is a basis of the topology of \(X\), and let \(\varphi : M \to N\) be a
  morphism of sheaves of \(\mathcal{O}_X\)-modules. If \(\varphi\) restricts to an
  isomorphism on the sections over every basic open \(B_i\), then \(\varphi\) is an
  isomorphism.
\end{lemma}

\begin{proof}
  \uses{lem:modules_isIso_iff_stalk}
  \leanok
  By Lemma~\ref{lem:modules_isIso_iff_stalk} it suffices to prove that the induced
  stalk map is bijective at each point \(x \in X\). For injectivity, a germ that
  vanishes at \(x\) already vanishes on some open neighbourhood, which may be
  shrunk to a basic open \(B_i\); on \(B_i\) the section map \(\varphi_{B_i}\) is
  injective, so the original section is zero, giving injectivity of the stalk map
  from injectivity over the basis. For surjectivity, an arbitrary germ at \(x\) is
  represented by a section of \(N\) over some basic open \(B_i\) containing \(x\);
  since \(\varphi_{B_i}\) is surjective, this section lifts to a section of \(M\)
  over \(B_i\), whose germ at \(x\) maps to the given germ. Thus the stalk map is
  bijective, and \(\varphi\) is an isomorphism.
\end{proof}

\begin{lemma}\leanok
[Affine-open locality criterion for isomorphisms of \(\mathcal{O}_X\)-modules]
  \label{lem:modules_isIso_iff_affineOpens}
  \lean{AlgebraicGeometry.Modules.isIso_iff_isIso_app_affineOpens}
  \uses{lem:modules_isIso_of_isBasis}
  Let \(X\) be a scheme and \(\varphi : M \to N\) a morphism of sheaves of
  \(\mathcal{O}_X\)-modules. Then \(\varphi\) is an isomorphism if and only if it
  restricts to an isomorphism on the sections over every affine open of \(X\).
\end{lemma}

\begin{proof}
  \uses{lem:modules_isIso_of_isBasis}
  \leanok
  The forward implication is immediate, since the section functor over a fixed open
  preserves isomorphisms. For the converse, the affine opens of a scheme form a
  basis of its topology, so the claim is the special case of
  Lemma~\ref{lem:modules_isIso_of_isBasis} in which the basis is the family of
  affine opens.
\end{proof}

\subsection{The affine computation}

The whole result rests on the following elementary computation in the affine
setting, where the base-change map is visibly an isomorphism because tensoring is
associative. Before stating it we isolate the one ingredient not yet available as
a ready-made result on which the affine reduction turns: the description of the
pushforward of a tilde-module along a morphism of affine schemes.

We first record three bespoke \(\Gamma\)-fragment supports that compute the global
sections of an affine pushforward; they carry no external citation and stand on
their proof sketches.

\begin{lemma}

\leanok
[Naturality of the global-sections comparison square]
  \label{lem:globalSectionsIso_hom_comp_specMap_appTop}
  \lean{AlgebraicGeometry.globalSectionsIso_hom_comp_specMap_appTop}
  Let \(\varphi : R \to R'\) be a homomorphism of commutative rings, and write
  \(\mathrm{gs}_R : \Gamma(\operatorname{Spec} R, \top) \cong R\) and
  \(\mathrm{gs}_{R'} : \Gamma(\operatorname{Spec} R', \top) \cong R'\) for the
  global-sections isomorphisms of the structure sheaves. Then the ring square
  \[
    \mathrm{gs}_R \;\circ\; (\operatorname{Spec}\varphi)^{\sharp}_{\top}
    \;=\; \varphi \;\circ\; \mathrm{gs}_{R'}
  \]
  commutes, where \((\operatorname{Spec}\varphi)^{\sharp}_{\top}\) is the
  top-open component of the structure-sheaf comparison map underlying
  \(\operatorname{Spec}\varphi\).
\end{lemma}

\begin{proof}


\leanok
  This is the naturality of the unit/counit of the
  \(\Gamma \dashv \operatorname{Spec}\) adjunction read off at the top open: the
  global-sections map of \(\operatorname{Spec}\varphi\) is, under the
  global-sections isomorphisms on each side, conjugate to \(\varphi\) itself.
  Transposing the naturality square of the \(\operatorname{Spec}\) unit gives
  exactly the asserted ring equation.
\end{proof}

\begin{lemma}

\leanok
[Global sections of an affine pushforward]
  \label{lem:gammaPushforwardIso}
  \lean{AlgebraicGeometry.gammaPushforwardIso}
  \uses{lem:globalSectionsIso_hom_comp_specMap_appTop}
  Let \(\varphi : R \to R'\) be a homomorphism of commutative rings and let
  \(N\) be an \(\operatorname{Spec}(R')\)-module. Then there is a canonical
  isomorphism of \(R\)-modules
  \[
    \Gamma\bigl((\operatorname{Spec}\varphi)_* N\bigr)
    \;\cong\;
    \operatorname{restr}_\varphi\bigl(\Gamma N\bigr),
  \]
  identifying the \(R\)-module of global sections of the pushforward with the
  restriction of scalars along \(\varphi\) of the \(R'\)-module of global sections
  of \(N\).
\end{lemma}

\begin{proof}
  \uses{lem:globalSectionsIso_hom_comp_specMap_appTop}
  \leanok
  Because \((\operatorname{Spec}\varphi)^{-1}(\top) = \top\), the global sections of
  the pushforward \((\operatorname{Spec}\varphi)_* N\) are, on underlying abelian
  groups, the global sections of \(N\) with no transport along the open-set
  identification. The two sides therefore carry the same underlying abelian group;
  both unwind to nested restriction-of-scalars towers along the structure-sheaf
  global-sections maps. These towers are reconciled by applying the
  restriction-of-scalars composition identity twice and substituting the ring
  equation of
  Lemma~\ref{lem:globalSectionsIso_hom_comp_specMap_appTop}, which shows the
  resulting \(R\)-action is exactly restriction of scalars along \(\varphi\).
\end{proof}

\begin{lemma}

\leanok
[Global sections of an affine pushforward of a tilde-module]
  \label{lem:gammaPushforwardTildeIso}
  \lean{AlgebraicGeometry.gammaPushforwardTildeIso}
  \uses{lem:gammaPushforwardIso}
  Let \(\varphi : R \to R'\) be a homomorphism of commutative rings and let \(M\)
  be an \(R'\)-module. Then there is a canonical isomorphism of \(R\)-modules
  \[
    \Gamma\bigl((\operatorname{Spec}\varphi)_* \widetilde{M}\bigr)
    \;\cong\;
    \operatorname{restr}_\varphi M .
  \]
\end{lemma}

\begin{proof}
  \uses{lem:gammaPushforwardIso}
  \leanok
  Specialise Lemma~\ref{lem:gammaPushforwardIso} to \(N = \widetilde{M}\) and
  compose with the tilde-unit isomorphism
  \(\Gamma(\widetilde{M}) \cong M\) of \(R'\)-modules (which restricts to an
  isomorphism of \(R\)-modules after applying \(\operatorname{restr}_\varphi\)).
\end{proof}

\begin{lemma}

\leanok
[Sections of an affine pushforward over an arbitrary open]
  \label{lem:gammaPushforwardIsoAt}
  \lean{AlgebraicGeometry.gammaPushforwardIsoAt}
  \uses{lem:globalSectionsIso_hom_comp_specMap_appTop}
  Let \(\varphi : R \to R'\) be a homomorphism of commutative rings, let \(N\) be
  an \(\operatorname{Spec}(R')\)-module, and let \(U\) be an open of
  \(\operatorname{Spec}(R)\), with preimage
  \(V = (\operatorname{Spec}\varphi)^{-1}(U)\) in \(\operatorname{Spec}(R')\). Then
  there is a canonical isomorphism of \(R\)-modules
  \[
    \Gamma\bigl((\operatorname{Spec}\varphi)_* N,\, U\bigr)
    \;\cong\;
    \operatorname{restr}_\varphi\bigl(\Gamma(N,\, V)\bigr) .
  \]
  This is the open-indexed generalization of
  Lemma~\ref{lem:gammaPushforwardIso}, which is the case \(U = \top\) (so that
  \(V = \top\) as well). The family \(\{e_U\}_U\) of these isomorphisms is moreover
  natural in the open argument: the structure-sheaf restriction maps on the two
  sides commute with the isomorphisms, so the family is a natural isomorphism of the
  two presheaves of \(R\)-modules in the open variable (this naturality is
  established in the proof below). For \(U = D(a)\) it is exactly the linear
  equivalence \(e_{D(a)}\) of movement~(1) in the proof of
  Lemma~\ref{lem:pushforward_spec_tilde_iso}.
\end{lemma}

\begin{proof}
  \uses{lem:globalSectionsIso_hom_comp_specMap_appTop}
  \leanok
  The construction copies that of Lemma~\ref{lem:gammaPushforwardIso} verbatim,
  with the top open \(\top\) replaced throughout by the arbitrary open \(U\) (and
  its preimage \(V\)). The point that makes this replacement free is that the
  \(R\)-module structure on the sections is supplied \emph{uniformly} by
  restriction of scalars along the global-sections ring map at the top open,
  independently of the evaluation open; consequently both sides unwind to the same
  nested restriction-of-scalars towers, reconciled by applying the
  restriction-of-scalars composition identity twice and substituting the
  \(\top\)-level ring equation of
  Lemma~\ref{lem:globalSectionsIso_hom_comp_specMap_appTop}. Naturality in the open
  follows because every constituent of the construction is either a structure-sheaf
  restriction map or conjugation by a fixed ring map, and all of these commute with
  further restriction along an inclusion of opens.

  \medskip
  \emph{Naturality of the family \(\{e_U\}_U\) in the open (remark).} The family is
  in fact a natural isomorphism between two presheaves of \(R\)-modules on
  \(\operatorname{Spec}(R)\): the source presheaf
  \(U \mapsto \Gamma\bigl((\operatorname{Spec}\varphi)_* N,\, U\bigr)\), with the
  structure-sheaf restriction maps of the pushforward, and the target presheaf
  \(U \mapsto
    \operatorname{restr}_\varphi\bigl(\Gamma(N,\, (\operatorname{Spec}\varphi)^{-1}U)\bigr)\),
  with the structure-sheaf restriction maps of \(N\) transported through the preimage
  operation \((\operatorname{Spec}\varphi)^{-1}\) and read as \(R\)-linear maps by
  restriction of scalars along \(\varphi\). Concretely, for any inclusion of opens
  \(U' \subseteq U\), with \(V = (\operatorname{Spec}\varphi)^{-1}U\) and
  \(V' = (\operatorname{Spec}\varphi)^{-1}U'\), the square with horizontal maps
  \(e_U, e_{U'}\) and vertical maps the respective structure-sheaf restrictions
  commutes. The reason is structural: each \(e_U\) is a composite whose constituents
  are of exactly two kinds — structure-sheaf restriction maps (of \(N\),
  respectively of \((\operatorname{Spec}\varphi)_* N\)) and identity-on-carrier
  \(\operatorname{restrictScalars}\) repackagings (conjugation by the \emph{fixed},
  open-independent ring map supplied by the restriction-of-scalars reconciliation and
  by the \(\top\)-level ring equation of
  Lemma~\ref{lem:globalSectionsIso_hom_comp_specMap_appTop}). Restriction maps commute
  with further restriction by the presheaf functoriality axiom (using
  \(V' = (\operatorname{Spec}\varphi)^{-1}U' \subseteq
  (\operatorname{Spec}\varphi)^{-1}U = V\), as \((\operatorname{Spec}\varphi)^{-1}\) is
  monotone), and conjugation by a fixed ring map commutes with restriction because the
  ring map does not depend on \(U\); the naturality square is the paste of these
  component squares, and each \(e_U\) being an isomorphism makes the family a natural
  isomorphism. This open-naturality is exactly the compatibility invoked, at the
  inclusion \(D(a) \subseteq \top\), in the proof of
  Lemma~\ref{lem:pushforward_spec_tilde_iso}.
\end{proof}

\begin{lemma}

\leanok
[The tilde restriction from the top open to a basic open is a localization]
  \label{lem:tildeRestriction_isLocalizedModule}
  \lean{AlgebraicGeometry.tildeRestriction_isLocalizedModule}
  Let \(M\) be an \(R'\)-module and \(b \in R'\). Then the structure-sheaf
  restriction map
  \[
    \Gamma(\widetilde M,\, \top) \longrightarrow \Gamma(\widetilde M,\, D(b)),
  \]
  read as a map of \(R'\)-modules, exhibits \(\Gamma(\widetilde M, D(b))\) as the
  localization of \(\Gamma(\widetilde M, \top)\) at \(\mathrm{powers}(b)\). This is
  the \(R'\)-side localization fact transported in movement~(3) of the proof of
  Lemma~\ref{lem:pushforward_spec_tilde_iso}.
\end{lemma}

\begin{proof}

\leanok
  By the defining property of \(\widetilde{(-)}\), the tilde-localization map
  \(M \to \Gamma(\widetilde M, D(b))\) exhibits its target as the localization of
  \(M\) at \(\mathrm{powers}(b)\). The global-sections map
  \(M \to \Gamma(\widetilde M, \top)\) is the localization of \(M\) at the trivial
  submonoid \(\mathrm{powers}(1)\) — since \(D(1) = \top\) — and is therefore
  bijective. The triangle identity relating these two localization maps to the
  structure-sheaf restriction \(\Gamma(\widetilde M, \top) \to
  \Gamma(\widetilde M, D(b))\) then exhibits the restriction itself, post-composed
  with the inverse of the bijection \(M \cong \Gamma(\widetilde M, \top)\), as the
  localization of \(\Gamma(\widetilde M, \top)\) at \(\mathrm{powers}(b)\).
\end{proof}

We record three further bespoke supports that drive the basic-open verification of
the affine-pushforward isomorphism below. They carry no external citation and stand
on their proof sketches.

\begin{lemma}

\leanok
[Restriction of scalars of a localized module]
  \label{lem:powers_restrictScalars}
  \lean{AlgebraicGeometry.IsLocalizedModule.powers_restrictScalars}
  Let \(A\) be a commutative \(R\)-algebra, \(S \subseteq R\) a submonoid, and
  \(f : M \to N\) an \(A\)-linear map that exhibits \(N\) as the localization of
  \(M\) at the image submonoid
  \(\operatorname{algMap}_A(S) \subseteq A\). Then the underlying \(R\)-linear map
  \(\operatorname{restr}_R f\) exhibits \(N\) as the localization of \(M\) at \(S\)
  itself.
\end{lemma}

\begin{proof}


\leanok
  This is the exact converse of the standard fact that a localization of \(M\) at a
  submonoid \(S\) of \(R\) remains a localization after enlarging the base to an
  \(R\)-algebra \(A\) along the image submonoid \(\operatorname{algMap}_A(S)\). One
  checks the three defining conditions of a localized module directly. Each
  \(s \in S\) acts invertibly on \(N\) because its image
  \(\operatorname{algMap}_A(s)\) does and the two actions coincide by the scalar
  tower \(R \to A \to N\). Surjectivity of \(f\) up to denominators and the
  equality-witness condition transport along the map
  \(s \mapsto \operatorname{algMap}_A(s)\) from \(S\) onto
  \(\operatorname{algMap}_A(S)\), again using
  \(\operatorname{algMap}_A(s)\cdot{} = s\cdot{}\). In the application
  \(A = R'\), \(S = \mathrm{powers}(a)\), one has
  \(\operatorname{algMap}_{R'}(\mathrm{powers}(a)) = \mathrm{powers}(\varphi a)\), so
  this identifies the localization of \(\operatorname{restr}_\varphi M\) at \(a\)
  with the localization of \(M\) at \(\varphi a\); it is the ring-change core of the
  basic-open computation below.
\end{proof}

\begin{lemma}

\leanok
[Section-level counit isomorphism from a localization hypothesis]
  \label{lem:fromTildeGamma_app_isIso_of_localized}
  \lean{AlgebraicGeometry.fromTildeΓ_app_isIso_of_isLocalizedModule}
  Let \(N\) be an \(\operatorname{Spec}(R)\)-module and \(a \in R\). Suppose the
  structure-sheaf restriction map \(\Gamma(N, \top) \to \Gamma(N, D(a))\), read as
  an \(R\)-linear map of the global-section modules, exhibits \(\Gamma(N, D(a))\) as
  the localization of \(\Gamma(N, \top)\) at \(\mathrm{powers}(a)\). Then the
  tilde--\(\Gamma\) counit
  \(\mathrm{fromTilde}\Gamma\,N : \widetilde{\Gamma(N,\top)} \to N\) restricts to an
  isomorphism on the sections over the basic open \(D(a)\).
\end{lemma}

\begin{proof}


\leanok
  Over \(D(a)\) the section map \(L\) of the counit sits in the triangle identity
  \(L \circ j = \rho\), where
  \(j : \Gamma(N,\top) \to \Gamma\bigl(\widetilde{\Gamma(N,\top)}, D(a)\bigr)\) is
  the tilde-localization map and \(\rho : \Gamma(N,\top) \to \Gamma(N, D(a))\) is
  the structure-sheaf restriction. Both \(j\) and \(\rho\) exhibit their targets as
  the localization of \(\Gamma(N,\top)\) at \(\mathrm{powers}(a)\) — \(j\) by the
  defining property of \(\widetilde{(-)}\), and \(\rho\) by hypothesis. By the
  uniqueness of localized modules there is a unique \(R\)-linear isomorphism \(e\)
  with \(e \circ j = \rho\); the triangle identity then forces \(L = e\), so \(L\) is
  an isomorphism. (This is the per-basic-open input feeding the basis-local
  criterion Lemma~\ref{lem:modules_isIso_of_isBasis}.)
\end{proof}

\begin{lemma}

\leanok
[Affine pushforward of a tilde-module, conditional form]
  \label{lem:pushforward_spec_tilde_iso_conditional}
  \lean{AlgebraicGeometry.pushforward_spec_tilde_iso_of_isLocalizedModule}
  \uses{lem:modules_isIso_of_isBasis, lem:gammaPushforwardTildeIso,
  lem:fromTildeGamma_app_isIso_of_localized}
  Let \(\varphi : R \to R'\) be a homomorphism of commutative rings and \(M\) an
  \(R'\)-module. Suppose that for every \(a \in R\) the structure-sheaf restriction
  of the pushforward \(N := (\operatorname{Spec}\varphi)_*\widetilde M\) exhibits its
  sections over \(D(a)\) as the localization of its global sections at
  \(\mathrm{powers}(a)\). Then there is a canonical isomorphism of
  \(\operatorname{Spec}(R)\)-modules
  \[
    (\operatorname{Spec}\varphi)_*\widetilde M \;\cong\;
    \widetilde{\bigl(\operatorname{restr}_\varphi M\bigr)} .
  \]
\end{lemma}

\begin{proof}
  \leanok
  \uses{lem:modules_isIso_of_isBasis, lem:gammaPushforwardTildeIso, lem:fromTildeGamma_app_isIso_of_localized}
  Write \(N := (\operatorname{Spec}\varphi)_*\widetilde M\). Under the hypothesis,
  Lemma~\ref{lem:fromTildeGamma_app_isIso_of_localized} applies for each \(a \in R\),
  so the counit \(\mathrm{fromTilde}\Gamma\,N\) restricts to an isomorphism on the
  sections over every basic open \(D(a)\). Since the basic opens form a basis of
  \(\operatorname{Spec}(R)\), the basis-local criterion
  Lemma~\ref{lem:modules_isIso_of_isBasis} upgrades this to: the counit
  \(\mathrm{fromTilde}\Gamma\,N\) is an isomorphism, i.e.\ \(N\) lies in the
  essential image of \(\widetilde{(-)}\). Composing the inverse counit with the
  tilde of the global-sections comparison
  Lemma~\ref{lem:gammaPushforwardTildeIso} yields the asserted object isomorphism.
\end{proof}

\begin{lemma}\leanok
[Affine pushforward of a tilde-module]
  \label{lem:pushforward_spec_tilde_iso}
  \lean{AlgebraicGeometry.pushforward_spec_tilde_iso}
  \uses{lem:modules_isIso_of_isBasis, lem:gammaPushforwardTildeIso,
        lem:pushforward_spec_tilde_iso_conditional, lem:powers_restrictScalars,
        lem:gammaPushforwardIso, lem:gammaPushforwardIsoAt,
        lem:tildeRestriction_isLocalizedModule}
% NOTE: the carrier-wall obstruction to `hloc` is resolved (restriction-of-scalars /
% scalar-tower instances feed lem:powers_restrictScalars; lem:tildeRestriction_isLocalizedModule
% supplies the R'-side localization). The single remaining obligation is the open-naturality of
% the section comparison e_{D(a)}, established as a prose remark in the proof of
% lem:gammaPushforwardIsoAt; the proof sketch below derives each per-a hloc(a) from the top-level
% localization plus that naturality square for the inclusion D(a) ⊆ ⊤, so no section-level square
% is hand-proved per a.
  Let \(\varphi : R \to R'\) be a homomorphism of commutative rings and let \(M\)
  be an \(R'\)-module. Then there is a canonical isomorphism of
  \(\operatorname{Spec}(R)\)-modules
  \[
    \bigl(\operatorname{Spec}\varphi\bigr)_*\,\widetilde{M}
    \;\cong\;
    \widetilde{\bigl(\operatorname{restr}_\varphi M\bigr)},
  \]
  where \(\operatorname{Spec}\varphi : \operatorname{Spec}(R') \to
  \operatorname{Spec}(R)\) is the induced morphism of affine schemes,
  \(\widetilde{(-)}\) denotes the quasi-coherent sheaf associated to a module, and
  \(\operatorname{restr}_\varphi M\) is the \(R\)-module obtained from \(M\) by
  restriction of scalars along \(\varphi\). That is, pushing the quasi-coherent
  sheaf \(\widetilde{M}\) forward along \(\operatorname{Spec}\varphi\) is, up to
  this canonical identification, the tilde of the restriction of scalars of \(M\).

  This is the classical fact that the direct image of a quasi-coherent sheaf along
  a morphism of affine schemes is computed on modules by restriction of scalars; it
  is the affine companion of the pullback formula (the affine case of the
  pushforward analogue of Stacks ``$\widetilde{M}$ and pullback''). It is recorded
  here as bespoke project infrastructure because the affine-pushforward
  identification does not appear as a ready-made result in the existing
  quasi-coherent sheaf theory.
\end{lemma}

\begin{proof}
  \leanok
  \uses{lem:modules_isIso_of_isBasis, lem:gammaPushforwardTildeIso, lem:pushforward_spec_tilde_iso_conditional, lem:powers_restrictScalars, lem:gammaPushforwardIso, lem:gammaPushforwardIsoAt, lem:tildeRestriction_isLocalizedModule}
  We build the isomorphism \emph{directly on a basis of basic opens}. This is the
  non-circular route: checking the comparison map on basic opens yields the object
  isomorphism, and quasi-coherence of the pushforward is then a corollary, not a
  prerequisite. (One must not reverse this order: the global-sections comparison
  \(\Gamma\bigl((\operatorname{Spec}\varphi)_*\widetilde{M}\bigr) \cong
  \operatorname{restr}_\varphi M\) of
  Lemma~\ref{lem:gammaPushforwardTildeIso} alone would only upgrade to an object
  isomorphism once the pushforward is already known to be quasi-coherent, so
  quasi-coherence cannot be deduced from an object isomorphism built that way.)

  \emph{The comparison morphism.} There is a canonical morphism of
  \(\operatorname{Spec}(R)\)-modules
  \[
    \alpha : \widetilde{\bigl(\operatorname{restr}_\varphi M\bigr)}
    \longrightarrow \bigl(\operatorname{Spec}\varphi\bigr)_*\widetilde{M},
  \]
  namely the tilde-adjoint of the global-sections identification of
  Lemma~\ref{lem:gammaPushforwardTildeIso}. We show \(\alpha\) is an isomorphism.

  \medskip
  \emph{Reduction to basic opens.} The basic opens \(D(a)\), for \(a \in R\), form a
  basis of the topology of \(\operatorname{Spec}(R)\). By the basis-local criterion
  Lemma~\ref{lem:modules_isIso_of_isBasis} it suffices to check that \(\alpha\)
  restricts to an isomorphism on the sections over each basic open \(D(a)\).

  \medskip
  \emph{The computation over \(D(a)\).} On the source, the sections of
  \(\widetilde{(\operatorname{restr}_\varphi M)}\) over \(D(a)\) are the
  localization \(\bigl(\operatorname{restr}_\varphi M\bigr)[1/a]\) as an
  \(R[1/a]\)-module. On the target, since \((\operatorname{Spec}\varphi)^{-1}D(a) =
  D(\varphi a)\), the sections of \(\bigl(\operatorname{Spec}\varphi\bigr)_*
  \widetilde{M}\) over \(D(a)\) are the sections of \(\widetilde{M}\) over
  \(D(\varphi a)\), i.e.\ the localization \(M[1/\varphi a]\), regarded as an
  \(R[1/a]\)-module by restriction of scalars along the induced map
  \(R[1/a] \to R'[1/\varphi a]\). These two localizations agree: the element \(a\)
  acts on \(\operatorname{restr}_\varphi M\) through \(\varphi a\), so localizing
  \(\operatorname{restr}_\varphi M\) at the multiplicative set generated by \(a\) is
  the same module as localizing \(M\) at the multiplicative set generated by
  \(\varphi a\). Both are the universal localization of
  \(M\) inverting \(\varphi a\), so the comparison map \(\alpha\) is an isomorphism
  on the sections over \(D(a)\). As this holds for every basic open, \(\alpha\) is
  an isomorphism.

  \medskip
  \emph{The explicit formalization route over \(D(a)\).} The agreement of the two
  localizations just described is the entire content, but it must be realized
  through the global ring \(R\), because the \(R\)-action on the pushforward
  sections over \(D(a)\) is built from the global-sections module of
  \(N := (\operatorname{Spec}\varphi)_*\widetilde M\) rather than directly from
  \(R[1/a]\). The reduction therefore proceeds exactly as the conditional
  assembly Lemma~\ref{lem:pushforward_spec_tilde_iso_conditional} prescribes: it
  suffices to supply, for each \(a \in R\), the hypothesis \(\mathrm{hloc}(a)\) that
  the structure-sheaf restriction \(\Gamma(N, \top) \to \Gamma(N, D(a))\) exhibits
  \(\Gamma(N, D(a))\) as the localization of \(\Gamma(N, \top)\) at
  \(\mathrm{powers}(a)\). The \(R\)-action on \(\Gamma(N, D(a))\) is the honest
  restriction-of-scalars action carried by the scalar tower \(R \to R' \to {-}\)
  determined by \(\varphi\); against that structure the ring-change converse
  Lemma~\ref{lem:powers_restrictScalars} is applicable, so the carrier of the
  \(R\)-module is exactly the one to which the localization claim refers, and no
  recursively-defined or ad-hoc section-level scalar action intervenes.
  We obtain \(\mathrm{hloc}(a)\) by transporting the already-known \(R'\)-side
  localization (Lemma~\ref{lem:tildeRestriction_isLocalizedModule}), in
  element-free fashion, across the open-indexed section comparison of
  Lemma~\ref{lem:gammaPushforwardIsoAt} \emph{using its naturality in the open}
  (the naturality remark in the proof of
  Lemma~\ref{lem:gammaPushforwardIsoAt}). This proceeds in three
  movements.

  \emph{(1) The \(D(a)\)-level linear equivalence.} This is the component at
  \(U = D(a)\) of the open-indexed family of
  Lemma~\ref{lem:gammaPushforwardIsoAt}: the \(R\)-linear isomorphism
  \[
    e_{D(a)} : \Gamma\bigl((\operatorname{Spec}\varphi)_*\widetilde M,\, D(a)\bigr)
      \;\cong\;
      \operatorname{restr}_\varphi\bigl(\Gamma(\widetilde M,\, D(\varphi a))\bigr) ,
  \]
  using \((\operatorname{Spec}\varphi)^{-1}D(a) = D(\varphi a)\), which is
  definitional. No fresh \(D(a)\)-level construction is needed: this is precisely the
  open-indexed isomorphism already supplied by
  Lemma~\ref{lem:gammaPushforwardIsoAt} evaluated at the basic open \(D(a)\).

  \emph{(2) The naturality square for \(D(a) \subseteq \top\).} The compatibility
  that relates \(e_{D(a)}\) to the \(\top\)-level isomorphism
  \(e_\top\) (the global-sections comparison of
  Lemma~\ref{lem:gammaPushforwardIso}) is precisely the open-naturality square of the
  family \(\{e_U\}_U\) (established in the proof of
  Lemma~\ref{lem:gammaPushforwardIsoAt}), instantiated at the inclusion
  \(D(a) \subseteq \top\): the square
  \[
    \begin{array}{ccc}
      \Gamma\bigl((\operatorname{Spec}\varphi)_*\widetilde M,\, \top\bigr)
        & \xrightarrow{\;e_\top\;}
        & \operatorname{restr}_\varphi\bigl(\Gamma(\widetilde M,\, \top)\bigr) \\[4pt]
      {\scriptstyle \mathrm{res}_{\top \to D(a)}}\big\downarrow & &
        \big\downarrow{\scriptstyle \mathrm{res}_{\top \to D(\varphi a)}} \\[4pt]
      \Gamma\bigl((\operatorname{Spec}\varphi)_*\widetilde M,\, D(a)\bigr)
        & \xrightarrow{\;e_{D(a)}\;}
        & \operatorname{restr}_\varphi\bigl(\Gamma(\widetilde M,\, D(\varphi a))\bigr)
    \end{array}
  \]
  commutes, using \((\operatorname{Spec}\varphi)^{-1}D(a) = D(\varphi a)\), which is
  definitional. Because the family \(\{e_U\}_U\) is a natural isomorphism, this single
  square is available uniformly in \(a\); there is no per-\(a\) section-level identity
  to re-prove by hand. It intertwines the source restriction
  \(\mathrm{res}_{\top \to D(a)}\) on the pushforward with the (restriction-of-scalars
  reading of the) target restriction \(\mathrm{res}_{\top \to D(\varphi a)}\) on
  \(\widetilde M\).

  \emph{(3) Discharging \(\mathrm{hloc}(a)\) by transport.} The target restriction
  \(\mathrm{res}_{\top \to D(\varphi a)} : \Gamma(\widetilde M, \top) \to
  \Gamma(\widetilde M, D(\varphi a))\) exhibits \(\Gamma(\widetilde M, D(\varphi a)) =
  M[1/\varphi a]\) as the \(\mathrm{powers}(\varphi a)\)-localization of
  \(M = \Gamma(\widetilde M, \top)\) over \(R'\); this is exactly
  Lemma~\ref{lem:tildeRestriction_isLocalizedModule}. Apply the ring-change converse
  Lemma~\ref{lem:powers_restrictScalars} (with \(R\)-algebra \(A = R'\) and submonoid
  \(\mathrm{powers}(a)\), so that
  \(\operatorname{algMap}_{R'}(\mathrm{powers}(a)) = \mathrm{powers}(\varphi a)\)) —
  its \(\operatorname{Algebra} R\, R'\) and scalar-tower hypotheses being supplied by
  the restriction-of-scalars structure along \(\varphi\) — to read this same map as
  the \(\mathrm{powers}(a)\)-localization over \(R\). The commuting naturality square
  of movement (2) intertwines this target restriction with the source restriction
  \(\mathrm{res}_{\top \to D(a)} : \Gamma(N, \top) \to \Gamma(N, D(a))\) by way of the
  isomorphisms \(e_\top\) and \(e_{D(a)}\); transporting the localization property
  across that square — via the transport-along-an-equivalence principle for localized
  modules — yields that \(\mathrm{res}_{\top \to D(a)}\) is the
  \(\mathrm{powers}(a)\)-localization on \(\Gamma(N, D(a))\) over \(R\), i.e.\
  \(\mathrm{hloc}(a)\). No section-level computation specific to the individual \(a\)
  is performed: the only \(a\)-dependent input is the instantiation of the single
  natural-isomorphism square at \(D(a)\). Feeding the family
  \(\{\mathrm{hloc}(a)\}_{a \in R}\) to the conditional assembly
  Lemma~\ref{lem:pushforward_spec_tilde_iso_conditional} produces the unconditional
  object isomorphism \(\alpha\).

  \medskip
  \emph{Naturality in the open drives the transport.} The comparison assembled
  here, \((\operatorname{Spec}\varphi)_* \circ \widetilde{(-)} \cong
  \widetilde{(-)} \circ \operatorname{restr}_\varphi\), is natural in \emph{both}
  arguments — in the module \(M\) and in the open \(U\). The naturality in the
  module is the usual functoriality; the naturality in the open is the naturality
  remark in the proof of Lemma~\ref{lem:gammaPushforwardIsoAt}: the family
  \(\{e_U\}_U\) of Lemma~\ref{lem:gammaPushforwardIsoAt} commutes with the
  structure-sheaf restriction maps on both sides. It is precisely this
  open-naturality that drives the localization transport \emph{uniformly}: the
  per-\(a\) localization fact \(\mathrm{hloc}(a)\) follows from the \(\top\)-level
  localization Lemma~\ref{lem:tildeRestriction_isLocalizedModule} together with that
  naturality square for the inclusion
  \(D(a) \hookrightarrow \top\), rather than from a section-level naturality square
  re-proved by hand for each individual \(a\). This is the single remaining
  obligation of the construction; the carrier of the \(R\)-action on the pushforward
  sections is the restriction-of-scalars structure along \(\varphi\), against which
  Lemma~\ref{lem:powers_restrictScalars} applies directly.

  The \(R\)-module structure on the \(R'\)-side sections to which the ring-change
  converse Lemma~\ref{lem:powers_restrictScalars} is applied is the honest
  restriction-of-scalars structure: \(R\) acts through \(\varphi\), i.e.\ via the
  \(R\)-algebra structure on \(R'\) determined by \(\varphi\) and the associated
  scalar tower \(R \to R' \to M[1/\varphi a]\). It is \emph{not} an ad-hoc,
  recursively-defined scalar action. Lemma~\ref{lem:powers_restrictScalars} is
  stated against exactly this restriction-of-scalars / scalar-tower structure, and
  it is under that structure that localizing \(\operatorname{restr}_\varphi M\) at
  \(a\) coincides with localizing \(M\) at \(\varphi a\).

  % NOTE: This lemma is the affine case of the general fact that the pushforward of
  % a quasi-coherent sheaf along an affine morphism is quasi-coherent. In current
  % Mathlib master that fact is the IsLocalizing / fromTildeGamma characterization
  % (AlgebraicGeometry.isIso_fromTildeΓ_pushforward); it post-dates the project's
  % pinned Mathlib, so the project proves it in-tree by the basic-open route above.
  % The in-tree proof remains valid and is not made redundant by a future Mathlib
  % bump.
  \medskip
  \emph{Quasi-coherence as a corollary.} The object isomorphism just constructed
  exhibits \(\bigl(\operatorname{Spec}\varphi\bigr)_*\widetilde{M} \cong
  \widetilde{(\operatorname{restr}_\varphi M)}\) as the tilde of a module. Tilde
  modules are quasi-coherent and quasi-coherence of \(\operatorname{Spec}(R)\)-modules
  is closed under isomorphism, so the pushforward
  \(\bigl(\operatorname{Spec}\varphi\bigr)_*\widetilde{M}\) is quasi-coherent. This
  conclusion is stated \emph{after} the object isomorphism, of which it is a
  consequence — no separate ``pushforward preserves quasi-coherent'' input is
  needed for the affine lane.
\end{proof}

\begin{remark}
  This single isomorphism discharges both inputs that the affine reduction below
  requires of the pushforward of a tilde-module.
  \begin{enumerate}
    \item \emph{The global-sections comparison.} Applying the global-sections
      (module-of-sections) functor to the isomorphism identifies the
      \(R\)-module of sections of
      \(\bigl(\operatorname{Spec}\varphi\bigr)_*\widetilde{M}\) with
      \(\operatorname{restr}_\varphi M\); this is exactly the \(\Gamma\)-fragment
      comparison used to recognise the section-level base-change map.
    \item \emph{Quasi-coherence of the pushforward.} The tilde of any module is
      quasi-coherent, and quasi-coherence of \(\operatorname{Spec}(R)\)-modules is
      closed under isomorphism; hence
      \(\bigl(\operatorname{Spec}\varphi\bigr)_*\widetilde{M}\) is quasi-coherent.
      In particular no separate ``pushforward preserves quasi-coherent'' theorem is
      needed for the affine lane.
  \end{enumerate}
  Because both inputs the affine reduction asks of the pushforward are corollaries
  of this single object isomorphism — the section-level module identification and
  the quasi-coherence of the pushforward — this lemma is the sole bespoke
  ingredient the affine lane requires.
\end{remark}

\subsection{The pullback companion}

The affine reduction below requires, alongside the pushforward dictionary, the
\emph{pullback} companion: the affine description of the inverse image of a
tilde-module. It is the affine case of the classical fact that the pullback of a
quasi-coherent sheaf is computed on modules by base change of scalars.

The pullback companion is built by uniqueness of left adjoints, which is driven by
the following packaging of the per-object \(\Gamma\)-fragment comparison as a
natural isomorphism of right-adjoint functors. It is bespoke project
infrastructure and carries no external citation.

\begin{lemma}

\leanok
[The pushforward--global-sections natural isomorphism]
  \label{lem:gammaPushforwardNatIso}
  \lean{AlgebraicGeometry.gammaPushforwardNatIso}
  \uses{lem:gammaPushforwardIso}
  Let \(\varphi : R \to R'\) be a homomorphism of commutative rings. Then the
  per-object global-sections comparison of
  Lemma~\ref{lem:gammaPushforwardIso} assembles into a natural isomorphism of
  functors \((\operatorname{Spec} R').\mathrm{Modules} \to \operatorname{ModuleCat}(R)\)
  \[
    \bigl(\operatorname{Spec}\varphi\bigr)_* \;\circ\; \Gamma_R
    \;\cong\;
    \Gamma_{R'} \;\circ\; \operatorname{restr}_\varphi ,
  \]
  whose component at an \(\operatorname{Spec}(R')\)-module \(N\) is the isomorphism
  \(\Gamma\bigl((\operatorname{Spec}\varphi)_* N\bigr) \cong
  \operatorname{restr}_\varphi(\Gamma N)\) of
  Lemma~\ref{lem:gammaPushforwardIso}.
\end{lemma}

\begin{proof}
  \uses{lem:gammaPushforwardIso}
  \leanok
  The components are the isomorphisms supplied by
  Lemma~\ref{lem:gammaPushforwardIso}. Naturality in the module argument \(N\)
  holds on the nose: every constituent of the comparison
  (the restriction-of-scalars composition repackagings and the
  restriction-of-scalars congruence) is the identity on underlying elements,
  merely re-presenting the module structure on an unchanged carrier, so each
  naturality square is a pointwise identity. This is the right-adjoint natural
  isomorphism that drives the uniqueness-of-left-adjoints argument for the affine
  pullback dictionary Lemma~\ref{lem:pullback_spec_tilde_iso}: identifying the two
  right adjoints \((\operatorname{Spec}\varphi)_* \circ \Gamma_R\) and
  \(\Gamma_{R'} \circ \operatorname{restr}_\varphi\) by this isomorphism forces the
  corresponding left adjoints to be isomorphic.
\end{proof}

\begin{lemma}

\leanok
[Affine pullback of a tilde-module]
  \label{lem:pullback_spec_tilde_iso}
  \lean{AlgebraicGeometry.pullback_spec_tilde_iso}
  \uses{lem:pushforward_spec_tilde_iso}
  % SOURCE: [Stacks Project], Schemes, Lemma `lemma-widetilde-pullback', Tag 01I9
  % (\S\,Quasi-coherent sheaves on affines) (read from references/stacks-schemes.tex,
  % L1241--1256).
  % SOURCE QUOTE: "Let
  % $(X, \mathcal{O}_X) = (\Spec(S), \mathcal{O}_{\Spec(S)})$,
  % $(Y, \mathcal{O}_Y) = (\Spec(R), \mathcal{O}_{\Spec(R)})$
  % be affine schemes.
  % Let $\psi : (X, \mathcal{O}_X) \to (Y, \mathcal{O}_Y)$ be a
  % morphism of affine schemes, corresponding to the ring map
  % $\psi^\sharp : R \to S$ (see Lemma \ref{lemma-category-affine-schemes}).
  % \begin{enumerate}
  % \item We have $\psi^* \widetilde M = \widetilde{S \otimes_R M}$
  % functorially in the $R$-module $M$.
  % \item We have $\psi_* \widetilde N = \widetilde{N_R}$ functorially
  % in the $S$-module $N$.
  % \end{enumerate}"
  \textit{Source: Stacks Project, Schemes, Lemma ``\(\widetilde M\) and pullback''
  (Tag 01I9), part (1).}
  Let \(\varphi : R \to R'\) be a homomorphism of commutative rings and let \(M\)
  be an \(R\)-module. Then there is a canonical isomorphism of
  \(\operatorname{Spec}(R')\)-modules
  \[
    \bigl(\operatorname{Spec}\varphi\bigr)^{*}\,\widetilde{M}
    \;\cong\;
    \widetilde{\bigl(R' \otimes_R M\bigr)},
  \]
  where \(\operatorname{Spec}\varphi : \operatorname{Spec}(R') \to
  \operatorname{Spec}(R)\) is the induced morphism of affine schemes,
  \(\widetilde{(-)}\) denotes the quasi-coherent sheaf associated to a module, and
  \(R' \otimes_R M\) is the base change of \(M\) along \(\varphi\) — the
  \(R'\)-module obtained by extension of scalars. The isomorphism is natural in
  \(M\). That is, pulling the quasi-coherent sheaf \(\widetilde{M}\) back along
  \(\operatorname{Spec}\varphi\) is, up to this canonical identification, the tilde
  of the base change of \(M\). It is the pullback companion of the affine
  pushforward dictionary Lemma~\ref{lem:pushforward_spec_tilde_iso}, and is part~(1)
  of the same Stacks lemma of which the pushforward formula is part~(2).
\end{lemma}

% SOURCE QUOTE PROOF: "The first assertion follows from the identification in
% Lemma \ref{lemma-compare-constructions}
% and the result of Modules, Lemma \ref{modules-lemma-restrict-quasi-coherent}."
\begin{proof}
  \uses{lem:pushforward_spec_tilde_iso}
  \leanok
  Pullback along \(\operatorname{Spec}\varphi\) is, on quasi-coherent sheaves, the
  extension-of-scalars operation \(- \otimes_R R'\) on modules transported through
  the \(\widetilde{(-)}\) / \(\operatorname{moduleSpec}\Gamma\) equivalence between
  modules and quasi-coherent sheaves. Concretely, the global sections of
  \(\bigl(\operatorname{Spec}\varphi\bigr)^{*}\widetilde{M}\) are the base change
  \(R' \otimes_R M\): the pullback functor \((\operatorname{Spec}\varphi)^{*}\) is
  left adjoint to the pushforward \((\operatorname{Spec}\varphi)_{*}\), and under
  the global-sections identification the latter is restriction of scalars along
  \(\varphi\) (Lemma~\ref{lem:pushforward_spec_tilde_iso}); the left adjoint of
  restriction of scalars is extension of scalars \(M \mapsto R' \otimes_R M\). The
  comparison map is the unit of the base-change adjunction, and the universal
  property of base change identifies it with the canonical isomorphism
  \(\bigl(\operatorname{Spec}\varphi\bigr)^{*}\widetilde{M} \cong
  \widetilde{(R' \otimes_R M)}\). Naturality in \(M\) is the naturality of that
  unit. (One may equally read this off the standard-open description: on
  \(D(\varphi a)\) the localization \((R' \otimes_R M)[1/a]\) is
  \(R'[1/\varphi a] \otimes_{R[1/a]} M[1/a]\), which is the base change of the
  sections of \(\widetilde M\) over \(D(a)\).)
\end{proof}

\subsection{The affine base-change lemma and its remaining obligations}

The affine base-change lemma below reduces, via the locality criterion
Lemma~\ref{lem:modules_isIso_iff_affineOpens} and the two affine dictionaries
Lemmas~\ref{lem:pushforward_spec_tilde_iso} and~\ref{lem:pullback_spec_tilde_iso},
to two obligations that are isolated here as named blocks. Both are extracted from
the proof of Stacks ``Affine base change'', and both are absent from the pinned
Mathlib: the first is the base-change-of-the-base-change-map compatibility under
restriction to affine opens, the second is the coherence identification of the
abstract adjoint mate with the concrete cancellation isomorphism.

\begin{lemma}
[Affine-local compatibility of the base-change map]
  \label{lem:base_change_map_affine_local}
  \lean{AlgebraicGeometry.TODO.base_change_map_affine_local}
  \uses{def:pushforward_base_change_map, lem:modules_isIso_iff_affineOpens}
  % SOURCE: [Stacks Project], Cohomology of Schemes, Lemma
  % `lemma-affine-base-change' (\S\,Cohomology and base change, I), proof, the
  % "local on $S$ and $S'$" step (read from references/stacks-coherent.tex,
  % L920--926).
  % SOURCE QUOTE: "The statement is local on $S$ and $S'$. Hence we may
  % assume $X = \Spec(A)$, $S = \Spec(R)$,
  % $S' = \Spec(R')$ and $\mathcal{F} = \widetilde{M}$
  % for some $A$-module $M$."
  \textit{Source: Stacks Project, Cohomology of Schemes, Lemma ``Affine base
  change'' (proof, local-on-\(S\)-and-\(S'\) step).}
  Let \(f : X \to S\) be an affine morphism and \(\mathcal{F}\) a quasi-coherent
  \(\mathcal{O}_X\)-module, with a base-change square as above. The formation of
  the base-change map \texttt{pushforwardBaseChangeMap}
  (Definition~\ref{def:pushforward_base_change_map}) is compatible with restriction
  to affine opens of \(S'\): for each affine open \(U \subseteq S'\), the
  restriction of the base-change map to sections over \(U\) agrees with the
  base-change map of the square obtained by restricting along \(U\). Consequently,
  by the affine-open locality criterion
  Lemma~\ref{lem:modules_isIso_iff_affineOpens}, the assertion that the base-change
  map is an isomorphism over an arbitrary base-change square \((S, S', X, X')\)
  reduces to the affine--affine case \(S = \operatorname{Spec}(R)\),
  \(S' = \operatorname{Spec}(R')\), \(X = \operatorname{Spec}(A)\),
  \(\mathcal{F} = \widetilde{M}\) for an \(A\)-module \(M\).
\end{lemma}

\begin{proof}
  \uses{def:pushforward_base_change_map, lem:modules_isIso_iff_affineOpens}
  Being an isomorphism of \(\mathcal{O}_{S'}\)-modules is, by
  Lemma~\ref{lem:modules_isIso_iff_affineOpens}, a property that may be checked on
  sections over the affine opens of \(S'\). Restricting the cartesian square along
  an affine open \(U \subseteq S'\) produces a cartesian square over
  \(U\) whose total space is the corresponding open of \(X'\); since affineness of
  \(f\) is preserved by base change, the restricted morphism is again affine and
  \(\mathcal{F}\) is again a tilde-module over the affine chart of \(X\) lying over
  \(U\) (this is where quasi-coherence of \(\mathcal{F}\) is used). The
  base-change map is the adjoint mate of the unit of the
  \(((g')^*, (g')_*)\)-adjunction, and one must verify that forming this mate
  commutes with the restriction; granting that compatibility, the general statement
  reduces to the affine--affine case. The required compatibility of the abstract
  adjoint mate with open restriction is not packaged in the pinned Mathlib, and is
  itself a base-change-of-the-base-change-map naturality statement; it is recorded
  here as a named obligation.
\end{proof}

\begin{lemma}
[Adjoint-mate identification with the cancellation isomorphism]
  \label{lem:pushforward_base_change_mate_cancelBaseChange}
  \lean{AlgebraicGeometry.TODO.pushforward_base_change_mate_cancelBaseChange}
  \uses{def:pushforward_base_change_map, lem:pushforward_spec_tilde_iso,
  lem:pullback_spec_tilde_iso}
  % SOURCE: [Stacks Project], Cohomology of Schemes, Lemma
  % `lemma-affine-base-change' (\S\,Cohomology and base change, I), proof, the
  % "boils down to the equality" step (read from references/stacks-coherent.tex,
  % L927--938).
  % SOURCE QUOTE: "We use Schemes, Lemma \ref{schemes-lemma-widetilde-pullback}
  % to describe pullbacks and pushforwards of $\mathcal{F}$.
  % Namely, $X' = \Spec(R' \otimes_R A)$ and
  % $\mathcal{F}'$ is the quasi-coherent sheaf associated
  % to $(R' \otimes_R A) \otimes_A M$.
  % Thus we see that the lemma boils down to the
  % equality
  % $$
  % (R' \otimes_R A) \otimes_A M = R' \otimes_R M
  % $$
  % as $R'$-modules."
  \textit{Source: Stacks Project, Cohomology of Schemes, Lemma ``Affine base
  change'' (proof, ``boils down to the equality'' step).}
  In the affine--affine case \(S = \operatorname{Spec}(R)\),
  \(S' = \operatorname{Spec}(R')\), \(X = \operatorname{Spec}(A)\),
  \(\mathcal{F} = \widetilde{M}\), write \(\Gamma(\alpha)\) for the
  global-sections incarnation of the base-change map
  \texttt{pushforwardBaseChangeMap} (Definition~\ref{def:pushforward_base_change_map}).
  Transport \(\Gamma(\alpha)\) through the two affine dictionaries that identify its
  source and target: by the pullback dictionary
  Lemma~\ref{lem:pullback_spec_tilde_iso} together with the pushforward dictionary
  Lemma~\ref{lem:pushforward_spec_tilde_iso}, the source is identified with
  \((R' \otimes_R A) \otimes_A M\) and the target with \(R' \otimes_R M\), both as
  \(R'\)-modules. Under these identifications \(\Gamma(\alpha)\) equals the concrete
  cancellation isomorphism
  \[
    \operatorname{cancelBaseChange} :
    (R' \otimes_R A) \otimes_A M \;\xrightarrow{\ \sim\ }\; R' \otimes_R M,
    \qquad (r' \otimes a) \otimes m \;\mapsto\; r' \otimes (a \cdot m)
  \]
  (Mathlib's \texttt{TensorProduct.AlgebraTensorModule.cancelBaseChange}). Since
  \(\operatorname{cancelBaseChange}\) is an isomorphism with no flatness hypothesis,
  this identification closes the affine case.
\end{lemma}

\begin{proof}
  \uses{def:pushforward_base_change_map, lem:pushforward_spec_tilde_iso, lem:pullback_spec_tilde_iso}
  The base-change map is defined abstractly as an adjoint mate: it is the transpose,
  under the \((\text{pullback} \dashv \text{pushforward})\) adjunctions for \(g\) and
  \(g'\), of the morphism obtained by applying \(f_*\) to the unit
  \(\mathcal{F} \to (g')_*(g')^*\mathcal{F}\). In the affine--affine case the four
  dictionary isomorphisms (the pullback dictionary
  Lemma~\ref{lem:pullback_spec_tilde_iso} for \((g')^*\widetilde{M}\) and for
  \(g^*(f_*\mathcal{F})\), and the pushforward dictionary
  Lemma~\ref{lem:pushforward_spec_tilde_iso} for \(f_*\widetilde{M}\) and for
  \(f'_*\) of the base-changed module) identify the source and target of
  \(\Gamma(\alpha)\) with the iterated tensor products
  \((R' \otimes_R A) \otimes_A M\) and \(R' \otimes_R M\). The genuine remaining
  content is the coherence computation that unwinds the adjoint mate of the
  \(((g')^*, (g')_*)\)-unit through these four dictionary isos and shows that the
  resulting \(R'\)-linear map is exactly \(\operatorname{cancelBaseChange}\),
  \((r' \otimes a) \otimes m \mapsto r' \otimes (a \cdot m)\). This mate-versus-
  concrete-map identification has no ready counterpart in the pinned Mathlib; it is
  the precise content that a later Mathlib bridge identifying the
  pushforward base-change map with the tilde of \(\operatorname{cancelBaseChange}\)
  would supply. Granting it, the affine case closes because
  \(\operatorname{cancelBaseChange}\) is an isomorphism with no flatness.
\end{proof}

\begin{lemma}\leanok
[Affine base change]
  \label{lem:affine_base_change_pushforward}
  \lean{AlgebraicGeometry.affineBaseChange_pushforward_iso}
  \uses{def:pushforward_base_change_map, lem:base_change_map_affine_local,
        lem:modules_isIso_iff_affineOpens, lem:pullback_spec_tilde_iso,
        lem:pushforward_base_change_mate_cancelBaseChange,
        lem:pushforward_spec_tilde_iso}
  % SOURCE: [Stacks Project], Cohomology of Schemes, Lemma
  % `lemma-affine-base-change' (\S\,Cohomology and base change, I)
  % (read from references/stacks-coherent.tex, L906--918).
  % SOURCE QUOTE: "Let $f : X \to S$ be a morphism of schemes.
  % Let $\mathcal{F}$ be a quasi-coherent $\mathcal{O}_X$-module.
  % Assume $f$ is affine.
  % In this case $f_*\mathcal{F} \cong Rf_*\mathcal{F}$ is
  % a quasi-coherent sheaf, and for every base change diagram
  % (\ref{equation-base-change-diagram})
  % we have
  % $$
  % g^*f_*\mathcal{F} = f'_*(g')^*\mathcal{F}.
  % $$"
  \textit{Source: Stacks Project, Cohomology of Schemes, Lemma ``Affine base
  change''.}
  Assume \(f : X \to S\) is an affine morphism and \(\mathcal{F}\) is a
  quasi-coherent \(\mathcal{O}_X\)-module. Then for every base-change square as
  above the canonical map of
  Definition~\ref{def:pushforward_base_change_map} is an isomorphism
  \[
    g^*\bigl(f_*\mathcal{F}\bigr) \;\cong\; f'_*\bigl((g')^*\mathcal{F}\bigr).
  \]
  The quasi-coherence hypothesis on \(\mathcal{F}\) is essential and is not a mere
  convenience: the proof works over an affine open by writing
  \(\mathcal{F} = \widetilde{M}\) for a module \(M\), and a general
  \(\mathcal{O}_X\)-module need not be of this form
  \(\widetilde{M}\); without it neither the affine description of the pushforward
  nor the tensor-associativity computation below is available, and the conclusion
  fails.
\end{lemma}

% SOURCE QUOTE PROOF: "The vanishing of higher direct images is
% Lemma \ref{lemma-relative-affine-vanishing}.
% The statement is local on $S$ and $S'$. Hence we may
% assume $X = \Spec(A)$, $S = \Spec(R)$,
% $S' = \Spec(R')$ and $\mathcal{F} = \widetilde{M}$
% for some $A$-module $M$.
% We use Schemes, Lemma \ref{schemes-lemma-widetilde-pullback}
% to describe pullbacks and pushforwards of $\mathcal{F}$.
% Namely, $X' = \Spec(R' \otimes_R A)$ and
% $\mathcal{F}'$ is the quasi-coherent sheaf associated
% to $(R' \otimes_R A) \otimes_A M$.
% Thus we see that the lemma boils down to the
% equality
% $$
% (R' \otimes_R A) \otimes_A M = R' \otimes_R M
% $$
% as $R'$-modules."
% NOTE (updated iter-242): Lean proof is a documented sorry. The locality
% reduction (lem:modules_isIso_iff_affineOpens) lands, and BOTH affine
% dictionaries are now proved and axiom-clean: lem:pushforward_spec_tilde_iso
% (iter-241) and its pullback companion lem:pullback_spec_tilde_iso (iter-242).
% Those are no longer the blockers. TWO Mathlib-absent obligations remain, neither
% with a quick sub-step:
%   (1) AFFINE REDUCTION. S, S', X, X' are arbitrary (only f affine, F quasi-
%       coherent); reducing to S=Spec R, S'=Spec R', X=Spec A, F=M~ needs the
%       naturality / base-change compatibility of pushforwardBaseChangeMap under
%       restriction to affine opens of S', which is itself not packaged in Mathlib.
%   (2) ADJOINT-MATE <-> cancelBaseChange. Even once affine, the abstract adjoint-
%       mate pushforwardBaseChangeMap transported through the four dictionary isos
%       must be shown equal to tilde(TensorProduct.AlgebraTensorModule.cancelBaseChange)
%       -- a coherence computation unwinding the mate of the ((g')^*,(g')_*)-unit.
% cancelBaseChange itself is present in Mathlib with no flatness; the gap is the
% mate identification, not the closing algebra. See
% informal/affineBaseChange_pushforward_iso.md and the in-file sorry comment.
\begin{proof}
  \uses{def:pushforward_base_change_map, lem:modules_isIso_iff_affineOpens, lem:pushforward_spec_tilde_iso, lem:pullback_spec_tilde_iso, lem:base_change_map_affine_local, lem:pushforward_base_change_mate_cancelBaseChange}
  \emph{First reduction (locality).} By
  Lemma~\ref{lem:modules_isIso_iff_affineOpens} it suffices to prove that the
  base-change map of Definition~\ref{def:pushforward_base_change_map} is an
  isomorphism on sections over every affine open of \(S'\). Fix one such affine
  open. The assertion is then local on \(S\) and on \(S'\), so we may assume all
  schemes in sight are affine: \(X = \operatorname{Spec}(A)\),
  \(S = \operatorname{Spec}(R)\), \(S' = \operatorname{Spec}(R')\), with \(f\)
  corresponding to a ring map \(R \to A\) and \(\mathcal{F} = \widetilde{M}\) the
  quasi-coherent sheaf associated to an \(A\)-module \(M\). (It is here that
  quasi-coherence of \(\mathcal{F}\) is used: it is what licenses writing
  \(\mathcal{F} = \widetilde{M}\) over the affine open.) Computing the fibre product
  gives \(X' = \operatorname{Spec}(R' \otimes_R A)\), and the pullback
  \(\mathcal{F}' = (g')^*\mathcal{F}\) is the quasi-coherent sheaf associated to
  the \((R' \otimes_R A)\)-module \((R' \otimes_R A) \otimes_A M\). In this affine
  situation the base-change map is a morphism \(\alpha\) of quasi-coherent
  \((\operatorname{Spec} R')\)-modules.

  \medskip
  \emph{Reduction to a concrete module map via full-faithfulness of tilde.} The
  functor \(\widetilde{(-)} : \operatorname{ModuleCat}(R') \to
  (\operatorname{Spec} R').\mathrm{Modules}\) is fully faithful, with essential
  image precisely the quasi-coherent modules; equivalently, the counit comparison
  \(\widetilde{\Gamma(\mathcal{N})} \to \mathcal{N}\) is an isomorphism exactly when
  \(\mathcal{N}\) is quasi-coherent. Source and target of \(\alpha\) are
  quasi-coherent, so both counits are isomorphisms; by naturality of the counit, a
  square relates \(\alpha\) to \(\widetilde{\Gamma(\alpha)}\), and since the two
  vertical counit maps are isomorphisms,
  \[
    \mathrm{IsIso}\,\alpha
    \;\Longleftrightarrow\;
    \mathrm{IsIso}\bigl(\Gamma(\alpha)\bigr),
  \]
  where \(\Gamma(\alpha)\) is a concrete \(R'\)-linear map of
  \(\operatorname{ModuleCat}(R')\). This reframing keeps the whole computation at
  the level of modules and morphisms of affine schemes; in particular it requires
  no manipulation of section-level scalar-multiplication or module instances. The
  ring map underlying the pushforward at the top open is the global-sections map
  \(f_{\top}\) of the morphism (which for \(f = \operatorname{Spec}\varphi\) is
  conjugate to \(\varphi\) via the naturality of the \(\Gamma\)--\(\operatorname{Spec}\)
  unit), and the preimage of the top open under \(\operatorname{Spec}\varphi\) is
  again the top open, so the identification of sections needs no transport. Any
  residual scalar identity in \(\operatorname{ModuleCat}(R')\) is settled by the
  scalar-tower compatibility of restriction of scalars.

  \medskip
  \emph{Identification of the two sides via the affine dictionaries.} The two
  quasi-coherent sheaves whose comparison \(\alpha\) is the section-level
  base-change map are identified over the affine chart by the two affine
  dictionaries. On the one hand, the \emph{pullback dictionary}
  Lemma~\ref{lem:pullback_spec_tilde_iso} identifies the restriction of
  \((g')^*\mathcal{F} = (g')^*(\widetilde{M})\) to the chart with the tilde of the
  base-changed module: writing \(g' = \operatorname{Spec}(\theta)\) for the induced
  ring map \(\theta : A \to R' \otimes_R A\),
  \[
    (g')^*\widetilde{M} \;\cong\; \widetilde{\bigl((R' \otimes_R A) \otimes_A M\bigr)},
  \]
  the tilde of the \((R' \otimes_R A)\)-module \((R' \otimes_R A) \otimes_A M\);
  pushing forward along the affine \(f' = \operatorname{Spec}(R \to R' \otimes_R A)\)
  and then reading global sections recovers \((R' \otimes_R A) \otimes_A M\) as an
  \(R'\)-module by restriction of scalars
  (Lemma~\ref{lem:pushforward_spec_tilde_iso}). On the other hand, the
  \emph{pushforward dictionary} Lemma~\ref{lem:pushforward_spec_tilde_iso}
  identifies \(f_*\mathcal{F} = f_*(\widetilde{M})\) over \(S\) with
  \(\widetilde{\operatorname{restr}_{R \to A} M}\), and the pullback dictionary
  Lemma~\ref{lem:pullback_spec_tilde_iso} then identifies its pullback
  \(g^*(f_*\mathcal{F})\) with \(\widetilde{R' \otimes_R M}\); reading global
  sections gives the \(R'\)-module \(R' \otimes_R M\). Thus, under the dictionaries,
  the source and target of \(\Gamma(\alpha)\) are the two iterated tensor products
  \((R' \otimes_R A) \otimes_A M\) and \(R' \otimes_R M\), both viewed as
  \(R'\)-modules.

  \medskip
  \emph{Identification of \(\Gamma(\alpha)\) with the cancellation isomorphism.} The
  base-change map \(\Gamma(\alpha)\) is the global-sections incarnation of
  \texttt{pushforwardBaseChangeMap} (Definition~\ref{def:pushforward_base_change_map}),
  which is constructed abstractly as an \emph{adjoint mate}: it is the transpose,
  under the \((\text{pullback} \dashv \text{pushforward})\) adjunctions for \(g\) and
  \(g'\), of the morphism obtained by applying \(f_*\) to the unit
  \(\mathcal{F} \to (g')_*(g')^*\mathcal{F}\). The crux of the remaining Lean work is
  to match this abstractly-defined adjoint mate, transported through the two
  dictionaries above, with the \emph{concrete} canonical \(R'\)-linear comparison
  between the two iterated tensor products — that is, with the associativity /
  cancellation isomorphism
  \[
    \operatorname{cancelBaseChange} :
    (R' \otimes_R A) \otimes_A M \;\xrightarrow{\ \sim\ }\; R' \otimes_R M
  \]
  of \(R'\)-modules (Mathlib's
  \texttt{TensorProduct.AlgebraTensorModule.cancelBaseChange}). Concretely the map
  sends \((r' \otimes a) \otimes m \mapsto r' \otimes (a \cdot m)\), using the
  \(A\)-action on \(M\). Once the identification of the abstract adjoint mate with
  this concrete \(\operatorname{cancelBaseChange}\) is made — the sole nontrivial
  remaining obligation — the affine case closes immediately, because
  \(\operatorname{cancelBaseChange}\) is an isomorphism with no flatness hypothesis
  whatsoever; flatness of \(g\) enters only in the global (non-affine) statement
  via the {\v C}ech argument, not here.
\end{proof}

\section{Flat base change for the pushforward}

\begin{theorem}\leanok
[Flat base change, $i=0$ case]
  \label{thm:flat_base_change_pushforward}
  \lean{AlgebraicGeometry.flatBaseChange_pushforward_isIso}
  \uses{def:pushforward_base_change_map, lem:affine_base_change_pushforward}
  % SOURCE: [Stacks Project], Tag 02KH (Lemma "Flat base change", part for
  % $i = 0$) (read from references/stacks-coherent.tex, L947--970, and
  % references/stacks-coherent.md).
  % SOURCE QUOTE: "Consider a cartesian diagram of schemes
  % $$
  % \xymatrix{
  % X' \ar[d]_{f'} \ar[r]_{g'} & X \ar[d]^f \\
  % S' \ar[r]^g & S
  % }
  % $$
  % Let $\mathcal{F}$ be a quasi-coherent $\mathcal{O}_X$-module
  % with pullback $\mathcal{F}' = (g')^*\mathcal{F}$.
  % Assume that $g$ is flat and that $f$ is quasi-compact and quasi-separated.
  % For any $i \geq 0$
  % \begin{enumerate}
  % \item the base change map of
  % Cohomology, Lemma \ref{cohomology-lemma-base-change-map-flat-case}
  % is an isomorphism
  % $$
  % g^*R^if_*\mathcal{F} \longrightarrow R^if'_*\mathcal{F}',
  % $$
  % \item if $S = \Spec(A)$ and $S' = \Spec(B)$, then
  % $H^i(X, \mathcal{F}) \otimes_A B = H^i(X', \mathcal{F}')$.
  % \end{enumerate}"
  \textit{Source: Stacks Project, Tag 02KH (``Flat base change''), the
  \(i = 0\) case.}
  Consider the cartesian square
  \[
    \begin{array}{ccc}
      X' & \xrightarrow{\;g'\;} & X \\[2pt]
      {\scriptstyle f'}\big\downarrow & & \big\downarrow{\scriptstyle f} \\[2pt]
      S' & \xrightarrow{\;g\;} & S
    \end{array}
  \]
  with \(\mathcal{F}\) a quasi-coherent \(\mathcal{O}_X\)-module and
  \(\mathcal{F}' = (g')^*\mathcal{F}\). Assume that \(g\) is flat and that \(f\)
  is quasi-compact and quasi-separated. Then the base-change map of
  Definition~\ref{def:pushforward_base_change_map} is an isomorphism
  \[
    g^*\bigl(f_*\mathcal{F}\bigr) \;\xrightarrow{\;\sim\;}\;
    f'_*\bigl((g')^*\mathcal{F}\bigr) = f'_*\mathcal{F}' .
  \]
  Equivalently, in the affine situation \(S = \operatorname{Spec}(A)\),
  \(S' = \operatorname{Spec}(B)\) with \(A \to B\) flat, the comparison map
  \[
    H^0(X, \mathcal{F}) \otimes_A B \;\xrightarrow{\;\sim\;}\;
    H^0(X_B, \mathcal{F}_B)
  \]
  is an isomorphism, where \(X_B = X \times_{\operatorname{Spec}(A)}
  \operatorname{Spec}(B)\) and \(\mathcal{F}_B\) is the pullback of
  \(\mathcal{F}\).

  The hypothesis that \(\mathcal{F}\) is quasi-coherent is essential: it is what
  makes \(f_*\mathcal{F}\) quasi-coherent on the affine base and thus the tilde of
  a module, and it is what lets the affine reduction
  (Lemma~\ref{lem:affine_base_change_pushforward}) write \(\mathcal{F}\) as
  \(\widetilde{M}\) over each affine open. Without it the comparison map need not be
  an isomorphism.
\end{theorem}

% SOURCE QUOTE PROOF: "Having said this, part (1) follows from part (2). Namely,
% part (1) is local on $S'$ and hence we may assume $S$
% and $S'$ are affine. In other words, we have $S = \Spec(A)$
% and $S' = \Spec(B)$ as in (2).
% Then since $R^if_*\mathcal{F}$ is quasi-coherent
% (Lemma \ref{lemma-quasi-coherence-higher-direct-images}),
% it is the quasi-coherent $\mathcal{O}_S$-module associated to the
% $A$-module $H^0(S, R^if_*\mathcal{F}) = H^i(X, \mathcal{F})$
% (equality by
% Lemma \ref{lemma-quasi-coherence-higher-direct-images-application}).
% Similarly, $R^if'_*\mathcal{F}'$ is the quasi-coherent
% $\mathcal{O}_{S'}$-module associated to the $B$-module
% $H^i(X', \mathcal{F}')$. Since pullback by $g$ corresponds
% to $- \otimes_A B$ on modules
% (Schemes, Lemma \ref{schemes-lemma-widetilde-pullback})
% we see that it suffices to prove (2).
%
% Let $A \to B$ be a flat ring homomorphism.
% Let $X$ be a quasi-compact and quasi-separated scheme over $A$.
% Let $\mathcal{F}$ be a quasi-coherent $\mathcal{O}_X$-module.
% Set $X_B = X \times_{\Spec(A)} \Spec(B)$ and denote
% $\mathcal{F}_B$ the pullback of $\mathcal{F}$.
% We are trying to show that the map
% $$
% H^i(X, \mathcal{F}) \otimes_A B \longrightarrow H^i(X_B, \mathcal{F}_B)
% $$
% (given by the reference in the statement of the lemma)
% is an isomorphism.
%
% In case $X$ is separated, choose an affine open covering
% $\mathcal{U} : X = U_1 \cup \ldots \cup U_t$ and recall that
% $$
% \check{H}^p(\mathcal{U}, \mathcal{F}) = H^p(X, \mathcal{F}),
% $$
% see
% Lemma \ref{lemma-cech-cohomology-quasi-coherent}.
% If $\mathcal{U}_B : X_B = (U_1)_B \cup \ldots \cup (U_t)_B$ we obtain
% by base change, then it is still the case that each $(U_i)_B$ is affine
% and that $X_B$ is separated. Thus we obtain
% $$
% \check{H}^p(\mathcal{U}_B, \mathcal{F}_B) = H^p(X_B, \mathcal{F}_B).
% $$
% We have the following relation between the {\v C}ech complexes
% $$
% \check{\mathcal{C}}^\bullet(\mathcal{U}_B, \mathcal{F}_B) =
% \check{\mathcal{C}}^\bullet(\mathcal{U}, \mathcal{F}) \otimes_A B
% $$
% as follows from
% Lemma \ref{lemma-affine-base-change}.
% Since $A \to B$ is flat, the same thing remains true on taking cohomology."
% NOTE: Lean proof is a documented sorry (deferred). Its affine input is the
% affine lemma above, which in turn reduces to lem:pushforward_spec_tilde_iso (the
% affine-pushforward-of-tilde brick) plus the tilde full-faithfulness reframe. On
% top of that it needs Čech-cohomology / affine-cover infrastructure for
% SheafOfModules, which Mathlib does not provide. This remains a multi-lane engine
% build, not a single-iter target. See informal/affineBaseChange_pushforward_iso.md.
\begin{proof}
  \uses{def:pushforward_base_change_map, lem:affine_base_change_pushforward}
  \emph{Reduction to the affine target.} The statement is local on \(S'\), so we
  may assume \(S\) and \(S'\) are affine, say \(S = \operatorname{Spec}(A)\) and
  \(S' = \operatorname{Spec}(B)\) with \(A \to B\) flat. Because \(f\) is
  quasi-compact and quasi-separated, \(f_*\mathcal{F}\) is quasi-coherent on
  \(S\); it is therefore the sheaf associated to the \(A\)-module
  \(H^0(S, f_*\mathcal{F}) = H^0(X, \mathcal{F})\). Likewise
  \(f'_*\mathcal{F}'\) is the sheaf on \(S'\) associated to the \(B\)-module
  \(H^0(X', \mathcal{F}')\). Since pullback along \(g\) corresponds to
  \(- \otimes_A B\) on modules, the base-change map becomes the comparison map of
  \(B\)-modules
  \[
    H^0(X, \mathcal{F}) \otimes_A B \longrightarrow H^0(X_B, \mathcal{F}_B),
    \qquad X_B = X \times_{\operatorname{Spec}(A)} \operatorname{Spec}(B),
  \]
  and it suffices to prove that this map is an isomorphism.

  \medskip
  \emph{The separated case via {\v C}ech complexes.} Suppose first that \(X\) is
  separated over \(A\). Choose a finite affine open cover
  \(\mathcal{U} : X = U_1 \cup \dots \cup U_t\); by separatedness all finite
  intersections \(U_{i_0 \dots i_p}\) are again affine, so {\v C}ech cohomology
  computes sheaf cohomology:
  \(\check{H}^p(\mathcal{U}, \mathcal{F}) = H^p(X, \mathcal{F})\). Base changing
  the cover yields an affine open cover
  \(\mathcal{U}_B : X_B = (U_1)_B \cup \dots \cup (U_t)_B\) of the again-separated
  \(X_B\), with the same property
  \(\check{H}^p(\mathcal{U}_B, \mathcal{F}_B) = H^p(X_B, \mathcal{F}_B)\). The
  affine case (Lemma~\ref{lem:affine_base_change_pushforward}) identifies each
  term of the base-changed {\v C}ech complex with the original term tensored over
  \(A\) with \(B\); that is, there is an isomorphism of cochain complexes
  \[
    \check{\mathcal{C}}^\bullet(\mathcal{U}_B, \mathcal{F}_B) \;\cong\;
    \check{\mathcal{C}}^\bullet(\mathcal{U}, \mathcal{F}) \otimes_A B .
  \]
  Because \(A \to B\) is flat, \(- \otimes_A B\) is an exact functor and hence
  commutes with the formation of cohomology of a complex; taking \(H^0\)
  (degree-zero cohomology) gives the desired isomorphism
  \(H^0(X, \mathcal{F}) \otimes_A B \cong H^0(X_B, \mathcal{F}_B)\).

  \medskip
  \emph{The quasi-separated case.} In general \(X\) need only be quasi-separated.
  Choose a finite affine open cover \(\mathcal{U} : X = U_1 \cup \dots \cup U_t\).
  Now the intersections \(U_{i_0 \dots i_p}\) are quasi-compact and separated
  rather than affine, but they fall under the separated case already treated, so
  flat base change holds for each of them. One then compares the two
  {\v C}ech-to-cohomology spectral sequences with \(E_2\)-pages
  \(\check{H}^p(\mathcal{U}, \underline{H}^q(\mathcal{F}))\) and
  \(\check{H}^p(\mathcal{U}_B, \underline{H}^q(\mathcal{F}_B))\), converging to
  \(H^{p+q}(X, \mathcal{F})\) and \(H^{p+q}(X_B, \mathcal{F}_B)\) respectively.
  The separated case provides the term-by-term identification
  \(\check{H}^p(\mathcal{U}_B, \underline{H}^q(\mathcal{F}_B)) =
  \check{H}^p(\mathcal{U}, \underline{H}^q(\mathcal{F})) \otimes_A B\) on the
  \(E_2\)-pages, and flatness of \(A \to B\) propagates this isomorphism through
  the spectral sequences to the abutment. In particular, in degree \(0\) one
  obtains \(H^0(X, \mathcal{F}) \otimes_A B \cong H^0(X_B, \mathcal{F}_B)\), as
  required. (For the \(i = 0\) statement at hand, one may avoid the spectral
  sequence entirely and argue by Mayer--Vietoris induction on the number \(t\) of
  cover members, reducing each step to the separated case.)
\end{proof}
